\subsection{偏导数}

1.6.1. 设 \(f\) 在 \(U\) 中点 \(x_0\) 的邻域 \(E\) 上定义且取值于 \(F\)。设 \(X\) 是 \(E\) 的向量子空间，\(V\) 是 \(x\) 中满足 \(X\) 的点 \(x_0 + x \in U\) 的集合；对 \(g(x) = f(x_0 + x)\) 置 \(x \in V\)。若 \(f\) 在 0 处有导数，则称 \(X\) 在 \(x_0\) 处有关于 \(g\) 的偏导数；此导数记作 \(D_X f(x_0)\)；它是从 \(X\) 到 \(F\) 的连续线性映射。若 \(f\) 在 \(x_0\) 处可微，则 f 在 \(X\) 处有关于 \(x_0\) 的偏导数，且此偏导数是 \(Df(x_0)\) 对 \(X\) 的限制。

1.6.2. 假设 \(E\) 是有限个赋范向量空间 \(E_{i}\)（\(1 \leqslant i \leqslant n\)）的乘积，并将它们按规范方式识别为 \(E\) 的子空间；设 \(x_{0} = (x_{0}^{1}, \ldots, x_{0}^{n})\) 属于 \(E\)，\(U\) 是 \(x_{0}\) 中 \(E\) 的一个邻域；最后设 \(f\) 是从 \(U\) 到 \(F\) 的映射。若存在，则用 \(D_{i}f(x_{0})\) 表示在点 \(x_{0}^{i}\) 处的映射 \(z_{i} \mapsto f(x_{0}^{1}, \ldots, z_{i}, \ldots, x_{0}^{n})\) 的导数，该映射在 \(x_{0}^{i}\) 中 \(E_{i}\) 的某邻域上定义且取值于 \(F\)。这是 \(\mathcal{L}(E_{i}; F)\) 中的元素，称为 \(f\) 在 \(x_{0}\) 处的第 i 个偏导数。若 \(f\) 在 \(x_{0}\) 处可微，则 \(n\) 个偏导数均存在，并由下式确定 \(Df(x_{0})\)：

\[
\mathrm{D} f (x _ {0}). h = \sum_ {i = 1} ^ {n} \mathrm{D} _ {i} f (x _ {0}). h _ {i} \quad \text { for } h = (h _ {1}, \dots , h _ {n}) \text { in } E.
\]

1.6.3. 更具体地，令 \(E = K^{n}\)。若 f 在 \(x_{0}\) 处的偏导数存在，则用 \(\partial_{i}f(x_{0})\) 表示 F 中的元素 \(D_{i}f(x_{0}).1\)。常使用以下记号；假设已为 \(K^{n}\) 上的坐标函数选定记号，例如 \(u_{i}\) 表示从 \(K^{n}\) 到 K 的第 i 个投影。于是写作：

\[
\frac {\partial f}{\partial u _ {i}} \left(x _ {0}\right) \quad \text { or } \quad \left. \frac {\partial f}{\partial u _ {i}} \right| _ {x = x _ {0}}
\]

而不写 \(\partial_{i}f(x_{0})\)。

1.6.4. 令 \(E = K^{n}\)、\(F = K^{m}\)；设取值于 F 的函数 \(f = (f_{1}, \ldots, f_{m})\) 在 E 中点 \(x_{0}\) 处可微。偏导数 \(a_{ji} = \partial_{i} f_{j}(x_{0})\) 于是存在（它们是 K 中的元素）。由 \(a_{ji}\) 构成的 m 行 n 列矩阵（元素位于指标为 j 的行和指标为 i 的列）称为 f 在 \(x_{0}\) 处的雅可比矩阵；相对于这些空间的规范基，它是从 \(Df(x_{0})\) 到 \(K^{n}\) 的线性映射 \(K^{m}\) 的矩阵。

